% Gemeinsame Präambel der M4-Dokumente (übernommen aus thesis/main.tex, damit die Integration einfach bleibt).
% Wird aus einem Unterordner eingebunden; Pfade sind relativ zu diesem Unterordner.
% --- Sprache und Schrift -------------------------------------------------------
\usepackage[T1]{fontenc}
\usepackage[utf8]{inputenc}
\usepackage[ngerman]{babel}
\usepackage{lmodern}
\usepackage{microtype}
\usepackage[autostyle=true,german=quotes]{csquotes}

% --- Layout --------------------------------------------------------------------
\usepackage[left=3cm,right=2.5cm,top=2.5cm,bottom=2.5cm]{geometry}
\usepackage{setspace}
\onehalfspacing

% --- Inhalte -------------------------------------------------------------------
\usepackage{graphicx}
\graphicspath{{abbildungen/}{../../../../../thesis/figures/}}
\usepackage{amsmath,amssymb}
\usepackage{siunitx}
\sisetup{locale=DE}
\usepackage{booktabs}
\usepackage{tabularx}
\usepackage{listings}
\lstset{basicstyle=\ttfamily\small,breaklines=true,frame=single,language=Python,
        numbers=left,numberstyle=\tiny,showstringspaces=false}
\usepackage{xcolor}
\usepackage{acro}
\usepackage{tikz}
\usetikzlibrary{arrows.meta,positioning,automata,fit,shapes.geometric,calc}
% Ablaufdiagramme (Programmablaufplan)
\tikzset{
  pap/.style={node distance=6mm and 10mm, font=\small},
  pap term/.style={draw, rounded corners=3mm, fill=black!8, align=center, minimum height=7mm,
                   inner xsep=3mm},
  pap proc/.style={draw, fill=white, align=center, text width=#1, minimum height=8mm},
  pap proc/.default=64mm,
  pap dec/.style={draw, diamond, aspect=2.6, fill=black!3, align=center, inner sep=1pt},
  pap arr/.style={-{Stealth[length=2.2mm]}, thick},
  pap lbl/.style={font=\scriptsize, fill=white, inner sep=1pt},
  pap good/.style={pap term, fill=green!12},
  pap bad/.style={pap term, fill=red!10},
}
\usepackage{algorithm}
\usepackage{algpseudocode}
\floatname{algorithm}{Algorithmus}
\algrenewcommand\algorithmicrequire{\textbf{Eingabe:}}
\algrenewcommand\algorithmicensure{\textbf{Ausgabe:}}
\algrenewcommand\algorithmicif{\textbf{wenn}}
\algrenewcommand\algorithmicthen{\textbf{dann}}
\algrenewcommand\algorithmicelse{\textbf{sonst}}
\algrenewcommand\algorithmicfor{\textbf{für}}
\algrenewcommand\algorithmicforall{\textbf{für alle}}
\algrenewcommand\algorithmicdo{\textbf{tue}}
\algrenewcommand\algorithmicwhile{\textbf{solange}}
\algrenewcommand\algorithmicreturn{\textbf{gib zurück}}
\algrenewcommand\algorithmicfunction{\textbf{Funktion}}
\algrenewcommand\algorithmicprocedure{\textbf{Prozedur}}
\algrenewcommand\algorithmicend{\textbf{Ende}}
\algrenewcommand{\algorithmiccomment}[1]{\hfill$\triangleright$ #1}
\renewcommand{\listalgorithmname}{Algorithmenverzeichnis}
\usepackage[colorinlistoftodos,textsize=footnotesize]{todonotes}

% --- Literatur (IEEE) ----------------------------------------------------------
\usepackage[backend=biber,style=ieee,dashed=false]{biblatex}
\addbibresource{../../../../../thesis/references.bib}
\addbibresource{../m4.bib}

% --- Links (zuletzt laden) -----------------------------------------------------
\usepackage[hidelinks]{hyperref}
\usepackage[ngerman]{cleveref}
\crefname{algorithm}{Algorithmus}{Algorithmen}


% --- Abkürzungen: die der Gesamtarbeit plus M4-spezifische ---------------------
\input{../../../../../thesis/chapters/00_abkuerzungen}
\input{../abkuerzungen_m4}

% --- Zusätzliche TikZ-Bibliotheken für die M4-Grafiken ---------------------------
\usetikzlibrary{decorations.pathreplacing}

% --- Kleine Hilfen -------------------------------------------------------------
\newcommand{\code}[1]{\texttt{\detokenize{#1}}}
\newcommand{\quelle}[1]{\textit{(Quelle: \code{#1})}}
\DeclareMathOperator*{\argmin}{arg\,min}
\newcommand{\E}{\mathbb{E}}
\newcommand{\N}{\mathcal{N}}
